% Generated draft: requires mathematical and source review.
\documentclass[11pt]{article}
\usepackage[margin=25mm]{geometry}
\usepackage{fontspec}
\setmainfont{Latin Modern Roman}
\usepackage{amsmath,amssymb,mathtools,unicode-math}
\setmathfont{Latin Modern Math}
\usepackage{enumitem,needspace,xurl,hyperref}
\setlength{\parindent}{0pt}
\setlength{\parskip}{5pt}
\setlength{\emergencystretch}{3em}
\newcommand{\N}{\mathbb{N}}
\newcommand{\Z}{\mathbb{Z}}
\newcommand{\Q}{\mathbb{Q}}
\newcommand{\R}{\mathbb{R}}
\newcommand{\C}{\mathbb{C}}
\newcommand{\F}{\mathbb{F}}
\newcommand{\A}{\mathbb{A}}
\newcommand{\PP}{\mathbb{P}}
\newcommand{\E}{\mathbb{E}}
\newcommand{\eps}{\varepsilon}
\newcommand{\dd}{\,\mathrm{d}}
\begin{document}
\section*{UnsolvedMath ID: 20001656}
\textbf{Source problem code:} catalogue entry
\subsection{Definitions and mathematical statement}
Let \(CV_n\), \(n\geq2\), be volume-one marked metric-graph Outer space, with translation lengths \(\ell_X(g)\) and forward asymmetric Lipschitz metric
\[
 d_L(X,Y)=\log\sup_{1\ne g\in F_n}\frac{\ell_Y(g)}{\ell_X(g)}.
\]
Fix \(o\in CV_n\). Its forward and backward normalized functions are
\[
 h_Y^+(X)=d_L(X,Y)-d_L(o,Y),\qquad
 h_Y^-(X)=d_L(Y,X)-d_L(Y,o).
\]
The closure of the forward image in the usual horofunction topology is the forward horocompactification; its complement of \(CV_n\) is the forward boundary. A Busemann point is a limit obtained from a suitable normalized almost-geodesic, with the source's formula \(d_L(X,\sigma(t))-t\) applying to a forward unit-speed ray based at \(o\). The source asks for this closure and boundary, also raises an unspecified simplicial metric, and asks whether every boundary point is Busemann.

The aid reports Horbez's solution for the \emph{forward} metric: the compactification is the primitive compactification, obtained by restricting projective translation-length functions to primitive elements of \(F_n\). Here primitive means belonging to some free basis. Its boundary is the Culler--Morgan tree boundary modulo
\[
 T\sim T'\quad\Longleftrightarrow\quad
 \exists c>0\ \forall p\text{ primitive}:\ell_T(p)=c\ell_{T'}(p),
\]
equivalently the reported special-pull equivalence after scaling. A normalized representative function is
\[
 h_T^+(X)=\log\sup_{p\text{ primitive}}\frac{\ell_T(p)}{\ell_X(p)}
 -\log\sup_{p\text{ primitive}}\frac{\ell_T(p)}{\ell_o(p)}.
\]
The aid reports that Busemann boundary points are exactly those represented by trees with dense orbits. Thus not all forward boundary points are Busemann. The exact special-pull definition is not included in the record.

For a fixed marked rose with basis \(a_1,\ldots,a_n\), the aid reports that its closed length simplex
\(\{(y_i):y_i\geq0,\ \sum_i y_i=1\}\)
injects into the forward compactification. Every proper face represents a non-Busemann point. On a positive test rose \(X_x\), the forward functions have the explicit restriction
\[
 h_y^+(X_x)=\log\max_i(y_i/x_i)-\log\max_i(y_i/o_i),
\]
when \(o\) is the fixed positive rose. The limiting tree is understood through its translation-length function when some edges have zero limiting lengths; it is not an interior free action.

The report further distinguishes the opposite orientation globally in rank three. In a rose with basis \(a,b,c\), let
\[
 y_k=(1-\varepsilon_k-\delta_k,\varepsilon_k,\delta_k),\qquad
 \varepsilon_k,\delta_k>0,\quad\varepsilon_k,\delta_k\to0,\quad
 \varepsilon_k/\delta_k\to r\in(0,\infty).
\]
The forward limit is the same face point \((1,0,0)\) for every \(r\). The usual finite rose-candidate expression for the reported global backward limit can be written
\[
 F_r(X)=\max\left\{\frac{\ell_X(b)}r,\ \ell_X(c),\
 \frac{\ell_X(bc)}{r+1},\ \frac{\ell_X(bc^{-1})}{r+1}\right\},
 \qquad H_r^-(X)=\log F_r(X)-\log F_r(o).
\]
This makes the residual collapse-rate dependence explicit; the finite-candidate lemma and the underlying artifact's exact formula require independent comparison. The aid reports pairwise distinct global limits as \(r\) varies, not merely different formulas before normalization. For the sum-symmetrized metric
\(d_\Sigma(X,Y)=d_L(X,Y)+d_L(Y,X)\),
the normalized limit is \(h_{(1,0,0)}^++H_r^-\), so those limits are also pairwise distinct. No assertion is made here for \(r=0\), \(r=\infty\), or every other symmetrization.

The historical discussion attributes the simplicial-metric and all-Busemann questions to Thierry Coulbois, records an optimistic comparison by ``Cormac'' with Teichmüller space in the Lipschitz metric, mentions CAT(0) analogies, and warns that normed spaces can have non-Busemann points. These are source-time comparisons without complete hypotheses, not general theorems for every CAT(0), Teichmüller or normed metric space. The aid supplies no complete simplicial-metric or backward-boundary classification.
\subsection{Short English statement}
The aid reports the forward Lipschitz horoboundary as the primitive tree compactification, with dense-orbit trees exactly the Busemann points. Rose-face examples preserve a common forward limit while separating backward and sum-symmetrized limits by relative edge-collapse rates.
\subsection{Sources}
\begin{itemize}
\item Local registry: permanent UnsolvedMath ID \(20001656\), exact public code \texttt{catalogue entry}; \texttt{aim-geometric-group-theory-notes.json}, \texttt{notes[164]}.
\item AIM workshop ``The geometry of the outer automorphism group of a free group'', Problem 5.3: \url{https://aimath.org/WWN/outerauto/outerauto.pdf}. The extraction loses closure-bar distinctions and includes the page-6 editor header. The discussion names Thierry Coulbois and ``Cormac'' without supplying the latter's surname in this record, and compares CAT(0), Lipschitz Teichmüller and normed-space boundaries. The literature field is empty.
\item Research attempt 1: \texttt{partial\_result}, status at run \texttt{partially\_solved}, rigor label \texttt{proved}, novelty confidence low, classification \texttt{PARTIAL-PROGRESS}. It attributes the forward primitive compactification and Busemann classification to Horbez and reports the marked-rose face and orientation proposition as a candidate contribution. No full Horbez bibliography or finite-candidate theorem citation is supplied. The current registry status is partially_solved.
\end{itemize}
\end{document}
